\documentclass[10pt,a4paper]{amsart}
\usepackage[margin=19mm]{geometry}
\usepackage{amsmath,amssymb,mathtools}
\usepackage[scr=rsfs,cal=euler]{mathalfa}
\usepackage{xfrac}
\usepackage[unicode,hidelinks,bookmarks=false]{hyperref}
\newcommand{\ie}{i.e.\ }
\newcommand{\cf}{cf.\ }
\newcommand{\resp}{respectively\ }
\newcommand{\Subsection}{\subsection}
\providecommand{\nobreakdash}{\nobreak\hbox{-}}
\setlength{\emergencystretch}{2em}
\multlinegap0pt
\usepackage{xcolor}
\usepackage{amssymb}
\newtheorem{theorem}{Theorem}
\title{\Large Adaptive Fast–Slow Operator Splitting for Multiscale Biochemical Stochastic Dynamics}
\author{Yuming Zeng, Wei Xie, Keqi Wang}
\date{Source: arXiv:2604.00140v1 (CC BY 4.0)}
\begin{document}
\maketitle

\noindent\emph{Source and licence.} \Large Adaptive Fast–Slow Operator Splitting for Multiscale Biochemical Stochastic Dynamics. Version arXiv:2604.00140v1 (CC BY 4.0), distributed by arXiv under the Creative Commons Attribution 4.0 International licence, http://creativecommons.org/licenses/by/4.0/, as recorded in the arXiv licence metadata for this exact version and shown on the abstract page https://arxiv.org/abs/2604.00140v1. Full licence notice: \texttt{licenses/arxiv-2604.00140-CC-BY-4.0.txt}.\par\smallskip

\noindent\textit{Editorial scope.} Counted unit: the article's theorem lem:truncation (source0.tex lines 499--524). The article's complete original proof of it is delivered with it (source0.tex lines 526--554, one \texttt{proof} environment of 29 lines). No mathematics is written, repaired or re-derived: the statement and the proof are the article's own lines, in the article's own notation, and no retained line is substituted or re-flowed.  Retained uncounted as the source-local context the proof needs: lines 488--494.  Nothing was omitted from any retained span, so the package carries no omission record.  No retained line contains a citation, so the wrapper carries no bibliography and no bibliography entry.  No retained line contains a figure environment, an image-inclusion command or drawing code, so no image is delivered and the package declares no asset dependency.  Editorial formatting only, in addition: an AMS \texttt{amsart} preamble that restates the article's own declarations and macros used by the retained lines, namely the environments \texttt{theorem} on separate counters, together with the standard packages those lines need.  The retained environments print in this document as Theorem 1 in order of appearance, whereas the article prints the same items with the numbering of its own full document.  The complete native source of the article as distributed in the arXiv e-print for this exact version is bundled unchanged as \texttt{sources/197565/source0.tex}.
\begin{equation}
\label{formula:truncate}
    \hat{Y}_{\Delta t} = Y^{(1)}
= \Delta t\,X_0
+ \sum_{k\in\mathcal{I}_{\mathrm{slow}}} \Delta W_k\,X_k
+ \sum_{j\in\mathcal{I}_{\mathrm{fast}}} \Delta W_j\,X_j,
\end{equation}

\begin{theorem}[Kunita Truncation Error]
\label{lem:truncation}
%Let $S_n=\pmb{s}$ denote the state at time $t_n$, and 
Let
$
S_{n+1} = \exp(Y_{\Delta t})(S_n)$ and $
\widehat{S}_{n+1} = \exp(\widehat{Y}_{\Delta t})(S_n),$
where $Y_{\Delta t}$ is the exact Stratonovich logarithm of the stochastic flow
over a time step $\Delta t$ starting from state $S_n=\pmb{s}$, and $\widehat{Y}_{\Delta t}$ denotes its first-order
truncation defined in Equation~\eqref{formula:truncate}.
Then the local mean-square error satisfies
\begin{equation}
\label{eq:Kunita_truncation}
\mathbb{E}\!\left[
\|S_{n+1}-\widehat{S}_{n+1}\|^2
\mid S_n=\pmb{s}
\right]
=
\frac14\,\Delta t^2
\sum_{i<j}
\|[X_i,X_j](\pmb{s})\|^2
+
O(\Delta t^3),
\end{equation}
where $[X_i,X_j]$ denotes the Lie bracket of the vector fields.
\end{theorem}

\begin{proof}[Proof sketch]
Write $Y_{\Delta t}=Y^{(1)}+Y^{(2)}+Y^{(3)}+\ldots$ and
$\hat Y_{\Delta t}=Y^{(1)}$.
A Fréchet expansion of the exponential map yields
$S_{n+1}-\hat S_{n+1}
=\exp(Y_{\Delta t})(\pmb{s})-\exp(\hat Y_{\Delta t})(\pmb{s})
=Y^{(2)}(\pmb{s})+\mathcal R_{\mathrm{trunc}}(\pmb{s})$, where
$\mathbb E\|\mathcal R_{\mathrm{trunc}}\|^{2}=o(\Delta t^3)$.
Consequently,
$\mathbb E[\|S_{n+1}-\hat S_{n+1}\|^{2}\mid S_n=\pmb{s}]
=\mathbb E[\|Y^{(2)}(\pmb{s})\|^{2}\mid S_n=\pmb{s}]+o(\Delta t^{3})$.
In Itô form, the second‑order term can be written as
$Y^{(2)}(\pmb{s})
=\sum_{1\le i<j\le N_r}H^{(i,j)}[X_i,X_j](\pmb{s})
 +\sum_{i=1}^{N_r} H^{(0,i)}[X_0,X_i](\pmb{s})$,
% where 
% $H^{(i,j)}=\tfrac12(I_{(i,j)}-I_{(j,i)})$ and $H^{(0,i)}=I_{(0,i)}$ are
% centered Itô iterated integrals
% \textcolor{red}{(we need to define $I_{(i,j)}$)}.
where 
$H^{(i,j)}=\tfrac12(I_{(i,j)}-I_{(j,i)})$ and $H^{(0,i)}=I_{(0,i)}$ are
centered Itô iterated integrals with
$I_{(i,j)} = \int_t^{t+\Delta t} \int_t^{\tau} dW_{s,i}\, dW_{\tau,j}$ and $
I_{(0,i)} = \int_t^{t+\Delta t} \int_t^{\tau} ds\, dW_{\tau,i}.$ 
Mixed covariances contribute only at higher
order, so only diagonal terms remain.
Using $\mathbb E[(H^{(i,j)})^{2}]=\tfrac14(\Delta t)^2$ and
$\mathbb E[(H^{(0,i)})^{2}]=\tfrac13(\Delta t)^3$ completes the proof.
\end{proof}

\end{document}
